\documentclass[11pt, letterpaper]{article}
\usepackage[margin=1in]{geometry}
\usepackage[T1]{fontenc}
\usepackage{lmodern}
\usepackage{parskip}
\usepackage{microtype}
\usepackage{soul}
\usepackage{xcolor}

\usepackage{ifthen}
\usepackage{amssymb}
\usepackage{xcolor}

\newboolean{showcomments}
\setboolean{showcomments}{true}

% use \setboolean{showcomments}{false} anywhere in your document after
% importing this to disable the comments

\makeatletter
\newcommand{\mynote}[3]{%
  \ifthenelse{\boolean{showcomments}}{%
   \fbox{\bfseries\sffamily\scriptsize#1}%
   {\small$\blacktriangleright$\textsf{\emph{\color{#3}{#2}}}$\blacktriangleleft$}}%
  {%
   % these two lines ensure that there is no blank space inserted
   \@bsphack
   \@esphack
  }%
}
\makeatother

% One command per author:
\definecolor{asparagus}{rgb}{0.53, 0.66, 0.42}
%\newcommand{\jc}[1]{\mynote{Jeronimo}{#1}{asparagus}}



\newcommand{\SMG}[1]{\mynote{SMG}{#1}{red}}    %Girvin
\newcommand{\shraddha}[1]{\mynote{SS}{#1}{blue}}     %Singh







\begin{document}

\noindent
Editorial Board\\
Physical Review X\\

\vspace{12pt}

\noindent
Dear Dr.\ Grocholski and Dr.\ Schiulaz,

We are pleased to resubmit our manuscript for consideration in Physical Review X. We thank you and the referees for the careful reading of the previous version, and for the opportunity to present our work at the editorial meeting. The feedback has been genuinely useful, and in response, we have substantially rewritten the introductory parts of the paper to make it more accessible. We believe the revised manuscript is significantly stronger as a result.

We begin by highlighting the aspects of the previous review that affirmed the scientific quality of this work. The referee explicitly acknowledged that our results are \ul{``well-motivated, non-trivial extensions of the existing literature''} and \ul{``a strong and timely contribution that meets the bar for PRX.''} We have taken these remarks as confirmation that the scientific content is sound and have focused the revision on the presentation and breadth concerns raised, which we address in full below. 

\textbf{On breadth.} The previous referee felt the paper was not a definitive landmark for the field of quantum control. We respectfully disagree and have strengthened the manuscript accordingly to bolster our viewpoint. The work goes beyond a trivial extension of the quantum control literature: Quantum Signal Processing (QSP) with non-commuting control variables is a fundamentally different and more advanced problem, studied as open problems in previous QSP literature. The theoretical language of non-Abelian QSP --- extending QSP to non-commuting control variables --- has been discussed in prior work with pursuit of open problems and CV-DV control through numerical optimizations~\cite{nemeth2023variants,hastrup2021improved,hastrup2021measurement,eickbusch2022fast}, but no explicit non-Abelian composite pulse sequence theory had previously been constructed. Our work demonstrates that non-commutativity in quantum phase space underpins a practical analytical control framework --- and is not merely a theoretical curiosity. Concretely, the framework already delivers a result of immediate experimental consequence, not only a future direction: piecewise gate teleportation yields the first high-fidelity, error-corrected single- and two-qubit logical gates for finite-energy GKP codes --- completing the analytical framework for universal logical control that the beyond-break-even GKP experiments demonstrated error correction for but lacked --- and does so with no additional hardware. An enabling primitive that opens a concrete experimental program is exactly the landmark criterion we understand PRX to apply. PRX publishes work that opens new research directions; this paper opens one that is immediately experimentally relevant in current superconducting and trapped-ion laboratories. The introduction now explicitly identifies the experimental and theoretical communities that stand to benefit from non-Abelian QSP: the ability to synthesize nonlinear polynomials in non-commuting position and momentum operators directly enables Hamiltonian simulation of anharmonic oscillators, interacting bosonic models, spin-boson systems, and molecular dynamics on non-parabolic potential energy surfaces --- problems of direct interest across quantum chemistry, condensed matter, and high energy physics. Furthermore, just as QSP was generalized to quantum singular value transformation (QSVT), unifying essentially all known qubit-only quantum algorithms, we anticipate that a complete constructive theory of NA-QSP may give rise to a non-Abelian QSVT unifying hybrid CV-DV quantum algorithms --- an important open problem this paper explicitly motivates through a discussion on phase estimation in Sec.~VIB.

\textbf{On conciseness and coherence.} A recurring concern was that the paper presented several small results combined into one large paper, with insufficient connections between them. We have addressed this directly and substantially. The revised manuscript stands at \ul{24 pages} of main text, reduced from 30 pages in the previous submission, and the full paper including appendices has been reduced from 48 to \ul{35 pages} --- a reduction achieved \emph{not} by cutting scientific content but by tightening the presentation and removing redundancy. More importantly, we have restructured the paper so that the Gaussian Controlled Rotation (GCR) primitive that we introduce as a non-Abelian QSP sequence is made explicit as the unifying thread running through every section. Each section now opens with a sentence connecting what it is doing back to GCR and hybrid control through the NA-QSP framework, and transitions between sections build on what came before rather than resetting. The reader's guide in Section~IA explicitly maps the paper's structure as a single coherent argument: GCR is introduced and analyzed in Section~III, serves as the engine for state preparation in Section~IV, unifies GKP stabilization, readout, and gate teleportation in Section~V, and reveals the geometry of universal oscillator control in Section~VI. This is not a collection of results --- it is a single framework applied systematically across the natural hierarchy of CV-DV control problems.

\textbf{On accessibility.} The previous referee found the paper too technical for a general PRX readership. We have taken this seriously. The introduction has been substantially rewritten to build the physical motivation from the ground up, starting from the connection between classical NMR composite pulses and the hybrid CV-DV control problem, before introducing the non-Abelian QSP framework. Section~II now includes a pedagogical reader's guide identifying the primary audiences --- QSP theorists familiar with DV qubits, readers familiar with CV or bosonic quantum information, and experimentalists working on superconducting and trapped-ion platforms --- and explaining what each reader needs from each section. The interpretation of GCR as an NA-QSP sequence has been made clearer in Section~III for better intuition to QSP theorists. We have also carefully distinguished the notation for classical pulse errors, quantum fluctuation errors, and entanglement error bounds, which were previously conflated.

We note that this paper sits at the intersection of quantum signal processing theory, hybrid CV-DV quantum control, and bosonic error correction --- three communities with partially overlapping but distinct backgrounds. A paper that genuinely serves all three will necessarily contain more explanatory scaffolding than one written for a single specialist audience. We have been careful to ensure that this scaffolding does real work for each reader type rather than adding length without purpose, and we ask that the paper be evaluated with this scope in mind. The one-page summary accompanying this submission is designed to help time-pressed referees navigate the paper and identify the results most relevant to their expertise.

\textbf{Summary of specific changes.}
\begin{itemize}
    \item \textbf{Abstract and Introduction} completely rewritten to unify previously competing voices, lead with the existence-proof argument, and explicitly connect to multiple reader communities.
    \item \textbf{Simplified figures} We have made the figures more simplified to relay the message in the text. Several figures (5,6) have been mildly polished without any change in content. The illustration in Figure 1 has been simplified. We have added a Figure 2 to show GCR with an illustration. And, with this addition Figure 4 (previously, figure 3) has become smaller and succinct. 
    \item \textbf{Section~II} reorganized with labeled pedagogical paragraphs, corrected notation throughout, physical interpretation of all key quantities, and an explicit bridge into Section~III.
    \item \textbf{Section~III} We have moved BB1(GCR) discussion completely to the appendix as just an extension, in order to get to the main applications in Sec. IV faster.
    \item The rotationally symmetric state preparation and Fock state results have been moved from Section~IV (D) to Section~VI (A), improving the flow of the paper: GKP state preparation in Section~IV (C) now leads directly into universal control of GKP states in Section~V.
    \item \textbf{Section~VI} titled \textit{Future Directions and Outlook} contains concrete examples of NA-QSP beyond the main results --- quantum phase estimation and universal state preparation --- illustrating the scope of the framework in quantum control more broadly towards generalization to arbitrary states and quantum algorithms. We have also expanded the analysis and comparisons of phase estimation in Section VI to depict its true placement in the current literature.
    \item \textbf{Section~VII} contains a focused conclusion with a summary of the work, and discussions on the scope of using the current work to improve numerical optimization of CV-DV control pulses, an addition suggested in a previous referee round.
\end{itemize}

We hope the revised manuscript and accompanying one-page summary give the editors and referees the context needed to evaluate this work on its merits. We are confident that the results are ready for implementation in current laboratories today, and that the analytical framework we introduce will be built upon by the community in the years ahead.

\vspace{12pt}
\noindent
Sincerely,

\vspace{6pt}
\noindent
Shraddha Singh, Baptiste Royer, and Steven M. Girvin

\bibliographystyle{elsarticle-num} 
\bibliography{CV_QSP.bib}
\end{document}